% Gap statement
Existing approaches to roof segmentation for solar potential estimation divide into two
families, each with a structural limitation.
Methods that rasterize airborne point clouds into a Digital Surface Model (DSM) before
segmenting \citep{Jakubiec2013, Freitas2015, Bergamasco2018} discard the vertical structure
encoded in raw point density: overhangs, dormers, and flush abutments between adjacent roof
planes collapse into ambiguous intensity gradients, and segmentation errors propagate directly
into tilt and orientation estimates.
Methods that operate on the three-dimensional (3D) point cloud directly \citep{Awrangjeb2014,
Verma2006, Sampath2010} achieve finer geometric fidelity but depend on manually annotated
training sets or hand-tuned building-specific parameters that do not transfer across cities
without re-annotation.
Neither family scales readily to the urban extent at which solar policy operates: annotating
thousands of rooftops per city is prohibitive, and DSM-derived estimates carry systematic
orientation errors on complex roof forms that compound across an entire municipality.

% Contribution paragraph
We propose an unsupervised 3D geometric clustering pipeline that segments roof planes directly
from raw airborne Light Detection and Ranging (LiDAR) point clouds, requiring no manually
labelled training data.
Building points are extracted and partitioned into planar segments through iterative
normal-based clustering, and each segment is attributed with tilt, azimuth, and projected
area without any intermediate rasterization step.
We validate the method on the Actueel Hoogtebestand Nederland 3 (AHN3) dataset across three
Dutch cities --- Rotterdam, Amsterdam, and Utrecht --- covering 12\,000 rooftops, demonstrating
that city-to-city transfer requires no retraining and no parameter re-tuning.
The specific contributions are as follows:
%
\begin{enumerate}
  \item An unsupervised roof segmentation algorithm that operates on raw 3D point clouds and
        produces per-segment tilt and azimuth estimates without labelled supervision or DSM
        rasterization.
  \item A large-scale validation across three geographically and morphologically distinct Dutch
        cities, covering 12\,000 rooftops drawn from AHN3, establishing transferability across
        urban contexts.
  \item A publicly released processing pipeline enabling city-scale photovoltaic (PV) potential
        estimation from nationally available airborne laser scanning data.
\end{enumerate}
